% ============================================================
%  chapters/ch07-geometry-of-expressions.tex
% ============================================================

\chapter{Geometry of Expressions}

\chapterepigraph{%
  An expression is not a string.\\
  It is a shape in a space of shapes.
}{Flyxion, \textit{Semantic Infrastructure}}

\sidenote{Connects to\\App.\,E on\\hypercubes as\\expression spaces}

The previous two chapters treated Spherepop syntactically (Chapter~5)
and dynamically (Chapter~6). This chapter treats Spherepop
\emph{geometrically}: expressions as points in a metric space, evaluation
as motion through that space, and the structure of computation as the
topology of reachable expression sets.

\section{Expressions as Topological Objects}

A Spherepop expression is not merely a syntactic object — it is a
topological object with a natural metric. Two expressions are close
if they differ in few places; distant if they differ substantially.

\begin{definition}{Edit Distance}{edit-dist}
  The \emph{edit distance} $d(E_1, E_2)$ between two Spherepop
  expressions is the minimum number of atomic operations (insert a
  sphere, delete a sphere, change a label) needed to transform
  $E_1$ into $E_2$.
\end{definition}

The set of all Spherepop expressions over a fixed signature, equipped
with edit distance, forms a metric space $(\mathcal{E}, d)$. Evaluation
defines a trajectory through this space:
\[
  E_0 \xrightarrow{\text{pop}} E_1 \xrightarrow{\text{pop}} \cdots
  \xrightarrow{\text{pop}} E_n.
\]

\section{Nesting Depth as a Geometric Quantity}

Depth is the primary geometric invariant of a Spherepop expression.
It measures the maximum containment level — the ``height'' of the
expression in the nesting hierarchy.

\begin{definition}{Expression Depth Profile}{depth-profile}
  The \emph{depth profile} of an expression $E$ is the function
  $\delta_E : \mathrm{nodes}(E) \to \mathbb{N}$ assigning depth to
  each sphere node. The \emph{maximum depth} is
  $D(E) = \max_{s} \delta_E(s)$.
\end{definition}

\begin{proposition}{Depth Decreases Under Evaluation}{depth-decrease}
  If $E \to_\text{pop} E'$, then $D(E') \leq D(E)$, with equality only
  when the popped redex was not the deepest sphere in $E$.
\end{proposition}

The depth profile describes the ``shape'' of an expression geometrically:
a tall, narrow expression has high maximum depth and few wide branches;
a flat, wide expression has low maximum depth and many branches at each
level.

\section{Constraint Shells}

From the perspective of Chapters~1–2, a Spherepop expression is a
system of nested constraint shells.

\begin{definition}{Constraint Shell}{constraint-shell}
  The \emph{constraint shell} at depth $k$ in expression $E$ is the
  set of spheres at depth exactly $k$. The innermost shell (maximum
  depth $D(E)$) contains the redexes; the outermost shell (depth 0)
  contains the root.
\end{definition}

Each shell constrains the evaluation of spheres in all deeper shells:
the operator at depth $k$ cannot apply until all depth-$k+1$ spheres
in its content list have been evaluated. The constraint structure is
therefore hierarchical — a tower of constraints, one per depth level.

\begin{figure}[h]
  \centering
  \begin{tikzpicture}[scale=1.0]
    % diagrams/expression-shells.tex
% Concentric shells showing depth levels of 1+(3×(2^2))

  % Shell 0 (outermost): + operator
  \draw[RSVPdeep, very thick] (0,0) circle (2.8cm);
  \node[RSVPdeep, font=\footnotesize\scshape] at (-2.3, 2.3) {Shell 0};
  \node[RSVPdeep, font=\large\bfseries] at (-1.5, 0) {$+$};
  % Atom 1
  \draw[RSVPmid, thick, fill=RSVPlight] (-1.5,-0.0) circle (0.0pt);
  \draw[RSVPaccent!60, thick] (-1.5, 0) circle (0.45cm);
  \node[font=\small] at (-1.5, 0) {$1$};

  % Shell 1: × operator
  \draw[RSVPmid, thick] (0.6, 0) circle (1.9cm);
  \node[RSVPmid, font=\footnotesize\scshape] at (2.0, 1.6) {Shell 1};
  \node[RSVPmid, font=\large\bfseries] at (-0.4, 0) {$\times$};
  % Atom 3
  \draw[RSVPaccent!60, thick] (-0.3, 0) circle (0.40cm);
  \node[font=\small] at (-0.3, 0) {$3$};

  % Shell 2 (innermost): ^ operator
  \draw[RSVPaccent, thick] (1.4, 0) circle (1.0cm);
  \node[RSVPaccent, font=\footnotesize\scshape] at (2.0, 0.8) {Shell 2};
  \node[RSVPaccent!80, font=\bfseries] at (1.15, 0) {$\uparrow$};
  % Atom 2 (base)
  \draw[RSVPaccent!50, thin] (0.95,-0.0) circle (0.25cm);
  \node[font=\tiny] at (0.95,0) {$2$};
  % Atom 2 (exp)
  \draw[RSVPaccent!50, thin] (1.7, 0) circle (0.25cm);
  \node[font=\tiny] at (1.7,0) {$2$};

  % Evaluation direction arrows
  \draw[-{Stealth[length=5pt]}, RSVPaccent, thick, opacity=0.7]
    (1.4,-1.2) -- (0.6,-2.0)
    node[midway, below right, font=\tiny\itshape, RSVPaccent] {pop 1};
  \draw[-{Stealth[length=5pt]}, RSVPmid, thick, opacity=0.7]
    (0.3,-2.1) -- (-0.5,-2.6)
    node[midway, below, font=\tiny\itshape, RSVPmid] {pop 2};
  \draw[-{Stealth[length=5pt]}, RSVPdeep, thick, opacity=0.7]
    (-0.8,-2.6) -- (-1.5,-2.9)
    node[midway, below left, font=\tiny\itshape, RSVPdeep] {pop 3};

  \end{tikzpicture}
  \caption{Constraint shells in the Spherepop expression
    $1 + (3 \times (2 \uparrow 2))$. Shell 0 (outermost, depth 0)
    contains the $+$ operator. Shell 1 contains $\times$. Shell 2
    (innermost) contains $\uparrow$ and the atoms $2, 2$. Evaluation
    proceeds shell 2 $\to$ shell 1 $\to$ shell 0.}
  \label{fig:expression-shells}
\end{figure}

\section{Evaluation Paths}

An evaluation path is a sequence of expressions from an initial
expression to its terminal form. The geometry of evaluation paths
reveals the structure of the computation.

\begin{definition}{Evaluation Path Length}{path-length}
  The \emph{length} of an evaluation path from $E$ to $E'$ is the
  number of pop steps. The \emph{shortest evaluation path} minimizes
  the number of steps.
\end{definition}

\begin{proposition}{Shortest Path Bound}{shortest-path}
  The length of the shortest evaluation path from $E$ to its terminal
  form equals the number of non-atomic spheres in $E$.
\end{proposition}

\begin{proof}
  Each pop step eliminates exactly one non-atomic sphere (replacing it
  with an atomic value). Since each step reduces the non-atomic count
  by exactly 1, and no step can be skipped in a valid evaluation, the
  minimum path length equals the initial number of non-atomic spheres.
\end{proof}

\section{The Expression Manifold}

At a higher level, the set of all expressions of a given depth and
operator signature forms an \emph{expression manifold}: a structured
space in which the geometry of evaluation can be studied globally.

\begin{definition}{Expression Manifold}{expr-manifold}
  For a fixed operator signature $(\Sigma_L, \mathrm{ar})$ and
  maximum depth $D$, the \emph{expression manifold} $\mathcal{M}(D)$
  is the metric space of all well-formed expressions of depth $\leq D$,
  equipped with edit distance $d$.
\end{definition}

Evaluation defines a vector field on $\mathcal{M}(D)$: at each
expression $E$, the ``pop vector'' points in the direction of $E$'s
reduction (toward the terminal expression). The integral curves of
this vector field are the evaluation paths.

\begin{remark}
  This makes the connection to RSVP explicit. In RSVP, the vector field
  $\mathbf{v}$ drives the evolution of the scalar field $\Phi$. In the
  expression manifold, the pop vector drives the evolution of the
  expression. Spherepop evaluation \emph{is} a special case of RSVP
  dynamics restricted to the discrete expression manifold.
\end{remark}

\section{Visualization Techniques}

Several visualization techniques help make expression geometry tangible:

\begin{enumerate}
  \item \textbf{Depth coloring:} color each sphere by depth, from
        cool colors at depth 0 to warm colors at maximum depth.
        The resulting image shows the thermal structure of the
        computation, with the ``hot'' innermost regions evaluating first.

  \item \textbf{Evaluation animation:} animate the sequence of pop
        steps, showing each sphere collapse in turn. The animation
        makes the propagation of values outward through the expression
        directly visible.

  \item \textbf{Expression tree as landscape:} embed the expression
        tree in 3D with depth as height. The landscape metaphor makes
        evaluation a physical process: water flowing downhill, filling
        valleys, eventually settling at the global minimum (the terminal
        expression at depth 0).

  \item \textbf{Constraint hypercube projection:} encode each expression
        of depth $\leq D$ as a binary vector of length $|\mathcal{M}(D)|$
        (one bit per possible sub-sphere, indicating presence or absence).
        The expression manifold embeds in the hypercube $Q_n$ from
        Appendix~E, and evaluation corresponds to motion toward a
        designated corner.
\end{enumerate}

\begin{exercise}
  Draw the expression tree for $(2 + 3) \times (4 + 1)$ as a
  3D landscape with depth as height. Identify the ``valleys''
  (innermost spheres) and describe the evaluation process as
  water flowing to sea level.
\end{exercise}

\begin{exercise}
  Two expressions $E_1$ and $E_2$ have the same depth profile if
  $\delta_{E_1} = \delta_{E_2}$ up to relabeling. Define an
  equivalence relation $\sim_\text{shape}$ on expressions that
  identifies expressions with the same shape. What does $\mathcal{E}/{\sim_\text{shape}}$
  look like? How many shape classes are there for depth $\leq 2$
  with binary operators?
\end{exercise}

\begin{exercise}[title={Harder}]
  Define a notion of \emph{expression curvature} that measures how
  ``unbalanced'' an expression tree is (i.e., the degree to which
  the depth profile is skewed to one side). Show that a perfectly
  balanced binary expression tree of depth $D$ has curvature 0,
  and that a completely left-leaning tree (every left child is a
  subtree, every right child is an atom) has maximum curvature for
  its depth.
\end{exercise}
